Per a mesurar la concentració de radó d'una estança, s'agafa una mostra de 180~cm$^3$ d'aire. Es col·loca la mostra dins d'un detector que compta el nombre total de desintegracions $\alpha$. Considerem que totes les desintegracions provenen del radó 222, $^{222}_{86}\mathrm{Rn}$. S'efectua un mesurament de les desintegracions durant 10 minuts cada dia, 10 dies seguits, i sempre a la mateixa hora. Les dades obtingudes són les següents:

\begin{center}
\begin{tabular}{l*{10}{c}}
\toprule
$t$ (dies) & 0 & 1 & 2 & 3 & 4 & 5 & 6 & 7 & 8 & 9 \\
$N$ (nuclis desintegrats) & 130 & 108 & 91 & 77 & 65 & 55 & 45 & 38 & 33 & 30 \\
\bottomrule
\end{tabular}
\end{center}

\figura[0.55]{fig1.png}

\begin{apartat}{1,25}
Representeu, dins la quadrícula adjunta, els nuclis desintegrats en funció del temps. A partir de la gràfica determineu el període de semidesintegració. Tenint en compte que l'evolució temporal del nombre de nuclis desintegrats és la mateixa que la del nombre total de nuclis de radó a l'estança, calculeu la constant de desintegració.
\end{apartat}

\begin{apartat}{1,25}
Tenint en compte que un mesurament dura 10 minuts, quina activitat té la mostra de 180~cm$^3$ d'aire de l'estança el dia $t = 0$? L'Agència de Protecció Ambiental dels Estats Units (EPA) recomana no sobrepassar l'activitat de 4~pCi per litre d'aire. Segons aquest límit, és perillosa la concentració a $t = 0$ a l'estança d'on s'ha extret la mostra?
\end{apartat}

\dades{%
  $1\ \mathrm{Ci} = 3,70\times 10^{10}\ \mathrm{Bq}$.}
